\let\negmedspace\undefined
\let\negthickspace\undefined
\documentclass[journal]{IEEEtran}
\usepackage[a5paper, margin=10mm, onecolumn]{geometry}
%\usepackage{lmodern} % Ensure lmodern is loaded for pdflatex
\usepackage{tfrupee} % Include tfrupee package

\setlength{\headheight}{1cm} % Set the height of the header box
\setlength{\headsep}{0mm}     % Set the distance between the header box and the top of the text

\usepackage{gvv-book}
\usepackage{gvv}
\usepackage{cite}
\usepackage{amsmath,amssymb,amsfonts,amsthm}
\usepackage{algorithmic}
\usepackage{graphicx}
\usepackage{textcomp}
\usepackage{xcolor}
\usepackage{txfonts}
\usepackage{listings}
\usepackage{enumitem}
\usepackage{mathtools}
\usepackage{gensymb}
\usepackage{comment}
\usepackage[breaklinks=true]{hyperref}
\usepackage{tkz-euclide} 
\usepackage{listings}
% \usepackage{gvv}                                        
\def\inputGnumericTable{}                                 
\usepackage[latin1]{inputenc}                                
\usepackage{color}                                            
\usepackage{array}                                            
\usepackage{longtable}                                       
\usepackage{calc}                                             
\usepackage{multirow}                                         
\usepackage{hhline}                                           
\usepackage{ifthen}                                           
\usepackage{lscape}
\begin{document}
\bibliographystyle{IEEEtran}
\title{10.7.51}
\author{EE25BTECH11002 - Achat Parth Kalpesh }
{\let\newpage\relax\maketitle}
\renewcommand{\thefigure}{\theenumi}
\renewcommand{\thetable}{\theenumi}
\setlength{\intextsep}{10pt} % Space between text and floats
\numberwithin{equation}{enumi}
\numberwithin{figure}{enumi}
\renewcommand{\thetable}{\theenumi}
\parindent 0px



\textbf{Question:}\\
Find the coordinates of all the points $\vec{P}$ on the ellipse $\frac{x^2}{a^2}+\frac{y^2}{b^2}$=1, for which the area of the triangle $PON$ is maximum, where $\vec{O}$ denotes the origin and $\vec{N}$, the foot of the perpendicular from $\vec{O}$ to the tangent at $\vec{P}$\\
\textbf{Solution:}\\
The equation of the ellipse is given as
\begin{align}
\vec{x}^\top\vec{V}\vec{x} + 2\vec{u}^\top\vec{x} + f &= 0\\
\vec{V} = \myvec{\frac{1}{a^2} & 0\\0 & \frac{1}{b^2}},\vec{u} &= \myvec{0\\0} , f = -1\\
\vec{x}^\top\vec{V} \vec{x} &= 1
\end{align}
The equation of a tangent to $\vec{P}$ is 
\begin{align}
    \brak{\vec{V}\vec{P} + \vec{u}}^\top \vec{x} + \vec{u}^\top\vec{P} + f = 0\\
    \brak{\vec{V}\vec{P}}^\top\vec{x} + f = 0
\end{align}
As $\vec{N}$ is the foot of the perpendicular from $\vec{O}$ to the tangent at $\vec{P}$
    \begin{align}
    \vec{N} &= \vec{0} - \frac{\brak{\vec{VP}}^T \vec{0} + f}{\norm{\vec{VP}}} \, \vec{VP}\\
    \vec{N} &= \frac{\vec{VP}}{\norm{\vec{VP}}}
\end{align} 


The cross product of two 2D vectors $\vec{x}$ and 
$\vec{y}$  can be written in matrix form using the skew-symmetric matrix

\begin{align}
    \vec{J} = \myvec{0 & 1 \\ -1 & 0}.
\end{align}

Hence,
\begin{align}
    \vec{x} \times \vec{y} = \vec{x}^\top \vec{J} \vec{y}.
\end{align}

The area of the triangle $\triangle PON$ is given by
\begin{align}
    \text{Area}\brak{PON}
    = \frac{1}{2} \frac{\abs{\vec{P}^\top \vec{J} \vec{V} \vec{P}}}{\norm{\vec{V}\vec{P}}^2},
\end{align}

Let
\begin{align}
    \vec{Q} &= \vec{V}\vec{P} \\
    \vec{P} &= \vec{V}^{-1}\vec{Q}.
\end{align}

Then, the area becomes
\begin{align}
    \text{Area}\brak{PON}
    &= \frac{1}{2} \frac{\abs{\brak{\vec{V}^{-1}\vec{Q}}^\top \vec{J} \vec{Q}}}{\norm{\vec{Q}}^2} \\
    &= \frac{1}{2} \frac{\abs{\vec{Q}^\top \brak{\vec{V}^{-1}\vec{J}} \vec{Q}}}{\norm{\vec{Q}}^2}.
\end{align}

Since $\vec{Q}^\top \vec{M} \vec{Q}$ depends only on the symmetric part of $\vec{M}$, define
\begin{align}
    \vec{S} = \frac{\vec{V}^{-1}\vec{J} + \brak{\vec{V}^{-1}\vec{J}}^\top}{2}.
\end{align}
Hence,
\begin{align}
    \text{Area}\brak{\vec{Q}} = \frac{1}{2}\frac{\abs{\vec{Q}^\top \vec{S} \vec{Q}}}{\norm{\vec{Q}}^2}.
\end{align}

Now,
\begin{align}
    \vec{V}^{-1} = \myvec{a^2 & 0 \\ 0 & b^2}, \quad 
    \vec{J} = \myvec{0 & 1 \\ -1 & 0}.
\end{align}
Thus,
\begin{align}
    \vec{V}^{-1}\vec{J} = \myvec{0 & a^2 \\ -b^2 & 0}, \quad
    \brak{\vec{V}^{-1}\vec{J}}^\top = \myvec{0 & -b^2 \\ a^2 & 0}.
\end{align}
So,
\begin{align}
    \vec{S} 
    &= \frac{1}{2}\brak{\vec{V}^{-1}\vec{J} + \brak{\vec{V}^{-1}\vec{J}}^\top} \\
    &= \frac{1}{2}\myvec{0 & a^2 - b^2 \\ a^2 - b^2 & 0} \\
    &= \frac{a^2 - b^2}{2}\myvec{0 & 1 \\ 1 & 0}.
\end{align}

The eigenvalues and eigenvectors of 
$\myvec{0 & 1 \\ 1 & 0}$ are
\begin{align}
    \lambda_1 = 1, \ \vec{v}_1 = \myvec{1 \\ 1}, \quad
    \lambda_2 = -1, \ \vec{v}_2 = \myvec{1 \\ -1}.
\end{align}
Hence, the eigenvalues of $\vec{S}$ are $\pm \frac{a^2 - b^2}{2}$, and the maximum absolute eigenvalue is
\begin{align}
    \lambda_{\max} = \frac{\abs{a^2 - b^2}}{2}.
\end{align}

The maximum area is therefore
\begin{align}
    \text{Area}_{\max} = \frac{\abs{a^2 - b^2}}{4}
\end{align}

The corresponding eigenvector is
\begin{align}
    \vec{Q} \propto \myvec{1 \\ 1}.
\end{align}
Hence,
\begin{align}
    \vec{P} = \vec{V}^{-1}\vec{Q} \propto \myvec{a^2 \\ b^2}.
\end{align}

Imposing the ellipse condition $\vec{P}^\top \vec{V} \vec{P} = 1$,
\begin{align}
    1 = \lambda^2 \brak{a^2 + b^2}
    \implies \lambda = \pm \frac{1}{\sqrt{a^2 + b^2}}.
\end{align}

Thus, the required points are
\begin{align}
    \vec{P} = \pm \frac{1}{\sqrt{a^2 + b^2}}\myvec{a^2 \\ b^2}
\end{align}
and the maximum area of $\triangle PON$ is
\begin{align}
    \text{Area}_{\max} = \frac{\abs{a^2 - b^2}}{4}
\end{align}

\begin{figure}[h!]
    \centering
    \includegraphics[width=\columnwidth]{figs/figure_py.png}
    \caption{Ellipse and triangles PON where a=3,b=2}
    \label{fig:fig}
 \end{figure}

\end{document}